\section*{Capítulo \ref{ch:g7:states-of-matter} --- Los estados de la materia y sus cambios}

\begin{solution}{exo:g7:states-of-matter:1}
\omterm{def:g3:solids-liquids-gases:solid}{Sólido}: \omterm{def:g7:states-of-matter:molecule}{moléculas} apretadas en una \omterm{def:g3:first-electric-circuit:battery}{pila} ordenada, temblando en
su sitio. \omterm{def:g3:solids-liquids-gases:liquid}{Líquido}: \omterm{def:g7:states-of-matter:molecule}{moléculas} en contacto pero
desordenadas, deslizándose unas junto a otras. \omterm{def:g3:solids-liquids-gases:gas}{Gas}:
\omterm{def:g7:states-of-matter:molecule}{moléculas} muy separadas, volando deprisa en todas las direcciones.
\end{solution}

\begin{solution}{exo:g7:states-of-matter:2}
En el \omterm{def:g3:solids-liquids-gases:solid}{sólido}, cada \omterm{def:g7:states-of-matter:molecule}{molécula} aguanta en su puesto dentro de la
\omterm{def:g3:first-electric-circuit:battery}{pila}, así que el conjunto conserva su forma. En el \omterm{def:g3:solids-liquids-gases:liquid}{líquido}, las
\omterm{def:g7:states-of-matter:molecule}{moléculas} se deslizan libremente unas junto a otras sin dejar de
tocarse: la multitud se desploma hasta encajar en el recipiente, hasta
que nada puede bajar más --- la superficie horizontal.
\end{solution}

\begin{solution}{exo:g7:states-of-matter:3}
El \omterm{def:g2:air-around-us:air}{aire} es casi todo espacio vacío entre \omterm{def:g7:states-of-matter:molecule}{moléculas} que vuelan:
empujar acerca a los voladores (y su martilleo responde --- de ahí la
elasticidad). Las \omterm{def:g7:states-of-matter:molecule}{moléculas} del agua ya se tocan; no queda vacío
que exprimir, así que el émbolo se encuentra con una pared.
\end{solution}

\begin{solution}{exo:g7:states-of-matter:4}
Que las \omterm{def:g7:states-of-matter:molecule}{moléculas} de \omterm{def:g3:solids-liquids-gases:gas}{gas} vuelan muy separadas, con un gran vacío
entre ellas --- se pueden apretar habitaciones enteras. El agua, cuyas
\omterm{def:g7:states-of-matter:molecule}{moléculas} ya se tocan, es incompresible: ninguna bomba la aprieta
más.
\end{solution}

\begin{solution}{exo:g7:states-of-matter:5}
El \omterm{def:g5:heat-and-insulation:heat}{calor} hace temblar cada vez más fuerte a las \omterm{def:g7:states-of-matter:molecule}{moléculas}
apiladas; en el \omterm{def:g4:changes-of-state:melting-point}{punto de fusión} el temblor puede más que el abrazo
de la \omterm{def:g3:first-electric-circuit:battery}{pila}, las filas se rompen y las \omterm{def:g7:states-of-matter:molecule}{moléculas} empiezan a
deslizarse: el \omterm{def:g3:solids-liquids-gases:solid}{sólido} se ha convertido en \omterm{def:g3:solids-liquids-gases:liquid}{líquido}.
\end{solution}

\begin{solution}{exo:g7:states-of-matter:6}
Un cambio de estado se limita a recolocar las mismas
\omterm{def:g7:states-of-matter:molecule}{moléculas} --- no se crea ni se destruye ninguna, así que la
\omterm{def:g3:measuring-mass:mass}{masa} no puede cambiar. Perder \omterm{def:g3:measuring-mass:mass}{masa} exigiría que
desaparecieran \omterm{def:g7:states-of-matter:molecule}{moléculas} (o que se escaparan del recipiente --- la
falsa alarma de la botella destapada).
\end{solution}

\begin{solution}{exo:g7:states-of-matter:7}
La viveza del movimiento molecular. Cuando lo caliente toca lo frío,
las \omterm{def:g7:states-of-matter:molecule}{moléculas} animadas aporrean a las adormiladas y les entregan
animación choque a choque --- las animadas se calman y las
adormiladas se espabilan --- hasta que las dos multitudes comparten un
mismo ritmo: \omterm{def:g3:temperature-thermometers:temperature}{temperaturas} iguales.
\end{solution}

\begin{solution}{exo:g7:states-of-matter:8}
Las \omterm{def:g7:states-of-matter:molecule}{moléculas} del perfume vuelan, chocan con las \omterm{def:g7:states-of-matter:molecule}{moléculas}
del \omterm{def:g2:air-around-us:air}{aire}, rebotan y vagabundean --- zigzag a zigzag por la
habitación. Una brisa se lleva multitudes enteras en bloque: montarse
en el \omterm{def:g2:air-around-us:wind}{viento} gana a ir dando tumbos entre la multitud.
\end{solution}

\begin{solution}{exo:g7:states-of-matter:9}
Solo las \omterm{def:g7:states-of-matter:molecule}{moléculas} \emph{más rápidas} pueden arrancarse del abrazo
del \omterm{def:g3:solids-liquids-gases:liquid}{líquido}, así que la evaporación es un impuesto sobre la viveza:
las fugitivas se llevan más de lo que les corresponde y la multitud
que queda atrás es, de media, más lenta --- es decir, más fría. El
sudor y el dedo mojado se enfrían exportando a sus más vivas.
\end{solution}

\begin{solution}{exo:g7:states-of-matter:10}
Al sol, las \omterm{def:g7:states-of-matter:molecule}{moléculas} atrapadas vuelan más deprisa y martillean la
piel con más \omterm{def:g2:pushes-and-pulls:force}{fuerza} y más a menudo: la granizada hincha el globo hasta
que la piel estirada la equilibra. Con frío, esas mismas
\omterm{def:g7:states-of-matter:molecule}{moléculas} se frenan, el martilleo se debilita y el empuje del
\omterm{def:g2:air-around-us:air}{aire} de fuera aplasta el globo hasta dejarlo más pequeño.
\end{solution}

\begin{solution}{exo:g7:states-of-matter:11}
Al \omterm{def:g2:ice-water-steam:meltfreeze}{congelarse}, el agua encaja en una \omterm{def:g3:first-electric-circuit:battery}{pila} insólitamente aireada y
abierta --- con más espacio vacío del que tenía la multitud
deslizante ---, así que el hielo ocupa más sitio con las mismas
\omterm{def:g7:states-of-matter:molecule}{moléculas}: \omterm{def:g6:density-floating:density}{densidad} por debajo de la del \omterm{def:g3:solids-liquids-gases:liquid}{líquido}, y
\omterm{def:g1:floating-and-sinking:words}{flota}. Si el agua fuera corriente, el hielo se apilaría apretado,
se hundiría según se formara y los estanques se congelarían de
abajo arriba --- sin dejarles a los peces ninguna agua resguardada.
\end{solution}

\begin{solution}{exo:g7:states-of-matter:12}
Nada removió el agua y, sin embargo, la tinta viajó: las propias
\omterm{def:g7:states-of-matter:molecule}{moléculas} del \omterm{def:g3:solids-liquids-gases:liquid}{líquido} están en movimiento incesante y empujan
a las \omterm{def:g7:states-of-matter:molecule}{moléculas} de tinta paso a paso, al azar, a través de la
multitud --- los dos \omterm{def:g3:solids-liquids-gases:liquid}{líquidos} están vivos de movimiento en cada
instante, aunque el vaso parezca perfectamente quieto. El \omterm{def:g5:heat-and-insulation:heat}{calor}
acelera todos los empujones, así que la expansión se apresura. La
tinta paciente hizo así visible la danza invisible --- la prueba que
convenció a quienes dudaban de las propias \omterm{def:g7:states-of-matter:molecule}{moléculas}.
\end{solution}

\begin{solution}{pb:g7:states-of-matter:1}
\textbf{1.} En filas ordenadas, hombro con hombro y cada uno moviendo
el cuerpo en el sitio --- y ninguno puede intercambiar el puesto ni
abandonarlo.
\textbf{2.} Las filas se deshacen: los alumnos siguen hombro con
hombro (¡tocándose!), pero ahora se entrelazan y se deslizan unos
junto a otros. Tienen que seguir manteniendo el contacto con la
multitud --- nadie puede echar a correr libre.
\textbf{3.} Los alumnos se dispersan, corriendo rápido y recto hasta
chocar entre sí o contra una pared, rebotando sin parar --- y, como
nada los mantiene juntos, solo las paredes los detienen: la multitud
reclama todos los rincones que ofrece el gimnasio.
\textbf{4.} Sí --- los mismos treinta en todas las escenas: la
\omterm{def:g3:measuring-mass:mass}{masa} sobrevive a los cambios de estado, representada.
\textbf{5.} La escena del \omterm{def:g3:solids-liquids-gases:gas}{gas} lo permite --- los corredores se
apiñan en la mitad del gimnasio (y martillean con más \omterm{def:g2:pushes-and-pulls:force}{fuerza} a quienes
los arrean). El apiñamiento del \omterm{def:g3:solids-liquids-gases:liquid}{líquido} no puede encoger: los
alumnos ya se tocan. Representado: los gases se aprietan, los
\omterm{def:g3:solids-liquids-gases:liquid}{líquidos} no.
\textbf{6.} Moverse cada vez más fuerte en su puesto; la escena se
convierte en la del \omterm{def:g3:solids-liquids-gases:liquid}{líquido} en el momento en que ese movimiento
revienta las filas y los alumnos empiezan a deslizarse unos junto a
otros.
\textbf{7.} La \omterm{def:g5:everyday-energy:energy}{energía} del mando se gasta en romper el abrazo de la
\omterm{def:g3:first-electric-circuit:battery}{pila} --- fila a fila --- y no en un movimiento más rápido: el crítico
\omterm{def:g3:temperature-thermometers:temperature}{termómetro} tiene que informar de que la viveza media \emph{no
cambia} hasta que se rompe la última fila. La meseta, representada.
\textbf{8.} Los alumnos \emph{más rápidos} --- solo ellos pueden
arrancarse. La viveza media del apiñamiento que queda baja: el
\omterm{def:g3:solids-liquids-gases:liquid}{líquido} se enfría a sí mismo al \omterm{def:g2:ice-water-steam:evapcond}{evaporarse}, exactamente el frío del
sudor y de los dedos mojados.
\textbf{9.} Los actores chocan una y otra vez con la cuerda y la
empujan hacia fuera; unos actores más rápidos empujan con más \omterm{def:g2:pushes-and-pulls:force}{fuerza} y
más a menudo, y el aro se hincha --- el globo al sol.
\textbf{10.} ``Apilaos en orden \emph{abierto} --- a un brazo de
distancia, aireados'': el \omterm{def:g3:solids-liquids-gases:solid}{sólido} del agua ocupa más sitio que su
\omterm{def:g3:solids-liquids-gases:liquid}{líquido}, así que la balsa de la escena de hielo flotaría sobre
la multitud de la escena líquida.
\textbf{11.} A la camiseta azul la van empujando --- golpe a golpe, al
azar, a través de la multitud deslizante --- hasta que ha vagabundeado
por todas partes: la tinta difundiéndose. Sube el mando y todos los
empujones se aceleran, así que el vagabundeo se apresura.
\textbf{12.} Por ejemplo: ``La obra de esta noche representa una sola
idea: toda la materia son \omterm{def:g7:states-of-matter:molecule}{moléculas} en movimiento. Concedida esa
idea, dos leyes viejas llegan gratis --- que la \omterm{def:g3:measuring-mass:mass}{masa} sobrevive a
todos los cambios de estado y la meseta de la \omterm{def:g6:water-states:changes}{fusión}. Lo que no
podemos representar es la escala: los actores de verdad son más
pequeños que cualquier mirada y más numerosos que todos los públicos
del mundo''.
\end{solution}
